% Level ceiling (grade 7): v = d/t as a quotient recipe (literal
% expressions arrived in math g7); average speed = total distance over
% total time; km/h <-> m/s only via gentle whole-number examples.
% The formal quotient-quantity treatment belongs to math grade 8.
\chapter{运动图像与平均速度}\label{ch:g7:motion-average-speed}

两年前，\omterm{def:g7:motion-average-speed:formula}{速度}指的是“一\omterm{ex:g2:measuring-time:clock}{小时}内走过的千米数”，
靠手指头数一数便可得出。如今，有了字母和图像作工具，
这一想法有了锋芒：一个能三向计算的公式，
一张图便能一眼看清整段行程——还有一个
关于平均值的著名陷阱，成年人也天天会上当。

\section{公式}

\begin{definition}[用公式表示速度]\label{def:g7:motion-average-speed:formula}
对于以匀速行进的一段路程（或其中一段），
\emph{速度}\index{速度} $v$ 等于路程 $d$ 除以
行驶时间 $t$：
\[
v = \frac{d}{t}.
\]
当 $d$ 以千米、$t$ 以\omterm{ex:g2:measuring-time:clock}{小时}为\omterm{def:g6:measurement-in-science:measuring}{单位}时，$v$ 的\omterm{def:g6:measurement-in-science:measuring}{单位}是
\emph{千米每\omterm{ex:g2:measuring-time:clock}{小时}}（\unit{km/h}）；若以米和
秒为\omterm{def:g6:measurement-in-science:measuring}{单位}，则为\emph{米每秒}（\unit{m/s}）——
科学家最青睐的\omterm{def:g6:measurement-in-science:measuring}{单位}。与\omterm{def:g7:volume-mass-density:formula}{密度}公式相似，这个公式可以
三向使用：$v = d/t$，\ $d = v \times t$，\ 
$t = d/v$。
\end{definition}

\begin{example}[公式的应用]\label{ex:g7:motion-average-speed:recipes}
一列火车 $2$ \omterm{ex:g2:measuring-time:clock}{小时}行驶了 \qty{240}{km}：$v = 240 \div 2 =
\qty{120}{km/h}$。以此\omterm{def:g7:motion-average-speed:formula}{速度}再行驶 $5$ \omterm{ex:g2:measuring-time:clock}{小时}，路程是多少？$d = 120
\times 5 = \qty{600}{km}$。以此\omterm{def:g7:motion-average-speed:formula}{速度}行驶 \qty{300}{km} 需要多长时间？$t = 300
\div 120 = 2.5$ \omterm{ex:g2:measuring-time:clock}{小时}——即两\omterm{ex:g2:measuring-time:clock}{小时}三十分钟（小心
这个捣蛋鬼：$0.5$ \omterm{ex:g2:measuring-time:clock}{小时}是 $30$ 分钟，绝不是“$50$”）。
\end{example}

\begin{example}[两种单位，同一速度]\label{ex:g7:motion-average-speed:units}
一名短跑运动员用 \qty{10}{s} 跑完 \qty{100}{m}：$v = \qty{10}{m/s}$。
用日常的\omterm{def:g6:measurement-in-science:measuring}{单位}表示，这个\omterm{def:g7:motion-average-speed:formula}{速度}是多少？一\omterm{ex:g2:measuring-time:clock}{小时}——即 $3600$ 秒——
内，运动员能跑 $10 \times 3600 = \qty{36000}{m} =
\qty{36}{km}$，所以 $\qty{10}{m/s} = \qty{36}{km/h}$。一般的换算关系为
每 \unit{m/s} 相当于 \qty{3.6}{km/h}——
因为一\omterm{ex:g2:measuring-time:clock}{小时}有 $3600$ 秒，而一千米只有
$1000$ 米。声速 \qty{340}{m/s}，换算成常用\omterm{def:g6:measurement-in-science:measuring}{单位}后超过
\qty{1200}{km/h}。
\end{example}

\section{用图像表示行程}

\begin{proposition}[图像规律]\label{prop:g7:motion-average-speed:graphs}
把路程随时间变化的情况画成图，运动就会“写下”
自己的经历：
\begin{enumerate}
  \item \omterm{def:g6:describing-motion:uniform}{匀速运动}画出一条\emph{直线}——
        相同时间内走过相同的路程，步步如此；
  \item 直线越\emph{陡}，运动就越快：
        陡峭程度就是\omterm{def:g7:motion-average-speed:formula}{速度}的直观呈现；
  \item 一段\emph{水平}线段表示停止：
        时间流逝，路程不变。
\end{enumerate}
拐弯处反映变化：向上弯曲表示加快；
趋于平缓则表示减慢。
\end{proposition}

\begin{omfigure}
\begin{tikzpicture}
\begin{axis}[omaxis, width=10.8cm, height=6.0cm,
    xlabel={时间 (\unit{min})}, ylabel={路程 (\unit{km})},
    xtick={0,10,20,30,40,50,60}, ytick={0,5,10,15,20},
    ymin=0, ymax=22, xmin=0, xmax=63]
  \addplot[very thick, omDef] coordinates
    {(0,0) (20,10) (35,10) (60,20)};
  \node[right, font=\footnotesize, omDef] at (axis cs:5,4.6)
    {骑行：陡峭};
  \node[above, font=\footnotesize, omDef] at (axis cs:27.5,10.2)
    {水平：停止};
  \node[right, font=\footnotesize, omDef] at (axis cs:44,13.6)
    {再次骑行：较缓};
\end{axis}
\end{tikzpicture}

{\small 一位快递员的早晨，由一条线道出：快速骑行，
十五分钟的停留，随后以较缓的\omterm{def:g7:motion-average-speed:formula}{速度}前行。这张图就是
这段行程的“自述”。}
\end{omfigure}

\begin{method}[读懂路程–时间图像]\label{met:g7:motion-average-speed:read}
\begin{enumerate}
  \item 先检查坐标轴及其\omterm{def:g6:measurement-in-science:measuring}{单位}——分钟还是
        \omterm{ex:g2:measuring-time:clock}{小时}，米还是千米；
  \item 在拐点处将折线分段，并为每段标注：
        直线且陡、直线且缓、水平；
  \item 对任一直线段，读出其纵向变化与横向变化，
        ——即路程的增量与时间的增量——并相除，得到
        该段的\omterm{def:g7:motion-average-speed:formula}{速度}；
  \item 对于整个过程，读出末时刻的总路程，
        并记住水平段也属于
        总时间的一部分。
\end{enumerate}
\end{method}

\begin{example}[解读快递员的行程]\label{ex:g7:motion-average-speed:decoded}
将该方法用于此图。第一段：\qty{10}{km}，用时
$20$ 分钟——即三分之一\omterm{ex:g2:measuring-time:clock}{小时}——于是 $v = 10 \div
\tfrac{1}{3} = \qty{30}{km/h}$。第二段为水平线段，从 $20$ 分钟
到 $35$ 分钟——这是送货停留。第三段：\qty{10}{km}，用时 $25$ 分钟，
\omterm{def:g7:motion-average-speed:formula}{速度}略低于 \qty{24}{km/h}。总计：\qty{20}{km}，用时一\omterm{ex:g2:measuring-time:clock}{小时}
——这恰好引出本章的下一个想法。
\end{example}

\section{平均速度——及其著名的陷阱}

\begin{definition}[平均速度]\label{def:g7:motion-average-speed:average}
整段行程的\emph{平均速度}\index{平均速度}
等于总路程除以\emph{总}时间——
停留时间也计入在内：
\[
v_{\text{平均}} = \frac{d_{\text{总}}}{t_{\text{总}}}.
\]
它是这样一个稳定的\omterm{def:g7:motion-average-speed:formula}{速度}：
以它匀速行进，在相同的总时间内即可走完同样的路程——快递员一\omterm{ex:g2:measuring-time:clock}{小时}内的 \qty{20}{km}，
即平均\omterm{def:g7:motion-average-speed:formula}{速度} \qty{20}{km/h}，尽管车轮从未
真正以这一\omterm{def:g7:motion-average-speed:formula}{速度}转动过。
\end{definition}

\begin{proposition}[陷阱]\label{prop:g7:motion-average-speed:trap}
一段行程的\omterm{def:g7:motion-average-speed:average}{平均速度}一般\emph{不}等于
其\omterm{def:g7:motion-average-speed:formula}{速度}的中点值。慢速路段比快速路段占用更多时间，
因此在平均值中权重更大。只有完整的算法——
即总路程除以总时间——才是可靠的；
直接对各\omterm{def:g7:motion-average-speed:formula}{速度}数值取平均，
是本章中最具诱惑力的错误做法。
\end{proposition}

\begin{example}[陷阱揭晓]\label{ex:g7:motion-average-speed:sprung}
一名骑行者去程以 \qty{30}{km/h} 的\omterm{def:g7:motion-average-speed:formula}{速度}骑行 \qty{60}{km}，
回程以 \qty{20}{km/h} 的\omterm{def:g7:motion-average-speed:formula}{速度}骑完同样的 \qty{60}{km}。有人说“\omterm{def:g7:motion-average-speed:average}{平均速度}：
\qty{25}{km/h}”？请认真核算。去程：$t = 60 \div 30 = 2$
\omterm{ex:g2:measuring-time:clock}{小时}。回程：$t = 60 \div 20 = 3$ \omterm{ex:g2:measuring-time:clock}{小时}。全程：
\qty{120}{km}，用时 $5$ \omterm{ex:g2:measuring-time:clock}{小时}——
$v_{\text{平均}} = \qty{24}{km/h}$。慢的那一半占去了
五\omterm{ex:g2:measuring-time:clock}{小时}中的三\omterm{ex:g2:measuring-time:clock}{小时}，把\omterm{def:g7:motion-average-speed:average}{平均速度}拉到了
中点值之下。在某一段上骑得越快，
这一段本身存在的时间就越短。
\end{example}

\begin{remark}[速度计说明了什么]\label{rem:g7:motion-average-speed:speedometer}
\omterm{def:g7:motion-average-speed:average}{平均速度}描述的是整段行程；\omterm{def:g7:motion-average-speed:formula}{速度}计的指针
回答的是另一个问题——\emph{此刻有多快}。在
图像上，“此刻”体现在图线于\emph{某一
点处}的陡峭程度上——在直线段上很容易看出，
在图线弯曲处则较为微妙。将“某点处的陡峭程度”精确化，
是数学最伟大的发明之一，它等着你
在高中阶段的末尾去学习。在此之前：直线段可以求出数值，
曲线只能靠文字描述。
\end{remark}

\section{练习}

\begin{exercise}[$\star$]\label{exo:g7:motion-average-speed:1}
写出\omterm{def:g7:motion-average-speed:formula}{速度}公式及其两个常用\omterm{def:g6:measurement-in-science:measuring}{单位}。
哪种方法可求距离？哪种可求时间？
\end{exercise}

\begin{exercise}[$\star$]\label{exo:g7:motion-average-speed:2}
一艘渡船 $3$ \omterm{ex:g2:measuring-time:clock}{小时}航行 \qty{45}{km}，一只野兔
跑 \qty{100}{m} 用时 \qty{8}{s}。计算两者的\omterm{def:g7:motion-average-speed:formula}{速度}，
各以其习惯\omterm{def:g6:measurement-in-science:measuring}{单位}表示。
\end{exercise}

\begin{exercise}[$\star$]\label{exo:g7:motion-average-speed:3}
利用换算关系，将 \qty{5}{m/s} 换算为 \unit{km/h}；\
将 \qty{72}{km/h} 换算为 \unit{m/s}。
\end{exercise}

\begin{exercise}[$\star$]\label{exo:g7:motion-average-speed:4}
叙述图像的三条规律。
图线逐渐变平的弯曲段说明了什么？
\end{exercise}

\begin{exercise}[$\star$]\label{exo:g7:motion-average-speed:5}
在快递员的图像上，哪两个时刻之间
\omterm{def:g7:motion-average-speed:formula}{速度}最小（但不为零）？如何不经计算即可判断？
\end{exercise}

\begin{exercise}[$\star$]\label{exo:g7:motion-average-speed:6}
某人步行的图像是一条直线，经过点
（$30$ 分钟，\qty{2}{km}）和（$60$ 分钟，\qty{4}{km}）。这是
匀速还是变速？\omterm{def:g7:motion-average-speed:formula}{速度}以 \unit{km/h} 计是多少？
\end{exercise}

\begin{exercise}[$\star$]\label{exo:g7:motion-average-speed:7}
给出\omterm{def:g7:motion-average-speed:average}{平均速度}的定义。一次徒步行程用时 $4$ \omterm{ex:g2:measuring-time:clock}{小时}，共走 \qty{12}{km}，
其中包括一\omterm{ex:g2:measuring-time:clock}{小时}的野餐。求\omterm{def:g7:motion-average-speed:average}{平均速度}，以及
\emph{步行时}的\omterm{def:g7:motion-average-speed:average}{平均速度}？
\end{exercise}

\begin{exercise}[$\star\star$]\label{exo:g7:motion-average-speed:8}
为什么长度相同时，慢的路段在平均值中的“权重”更大？
请用时间来回答，不要用公式。
\end{exercise}

\begin{exercise}[$\star\star$]\label{exo:g7:motion-average-speed:9}
重温这一陷阱：去程以 \qty{15}{km/h} 的\omterm{def:g7:motion-average-speed:formula}{速度}走 \qty{30}{km}，
回程以 \qty{30}{km/h} 的\omterm{def:g7:motion-average-speed:formula}{速度}走 \qty{30}{km}。预测中点值，
计算真实的\omterm{def:g7:motion-average-speed:average}{平均速度}——以及行程中哪一段
占用的时间更多？
\end{exercise}

\begin{exercise}[$\star\star$]\label{exo:g7:motion-average-speed:10}
画出（或精确描述）这次行程的图像：先
以 \qty{40}{km/h} 匀速行驶半\omterm{ex:g2:measuring-time:clock}{小时}；停留四分之一\omterm{ex:g2:measuring-time:clock}{小时}；
再以 \qty{60}{km/h} 匀速行驶四分之一\omterm{ex:g2:measuring-time:clock}{小时}。然后计算
这次行程的\omterm{def:g7:motion-average-speed:average}{平均速度}。
\end{exercise}

\begin{exercise}[$\star\star$]\label{exo:g7:motion-average-speed:11}
雷暴测距升级版：雷声在闪电后 \qty{6}{s} 才传来。
已知\omterm{def:g3:sound-around-us:sound}{声音}的\omterm{def:g7:motion-average-speed:formula}{速度}为 \qty{340}{m/s}，用 $d = v \times t$ 求雷暴的距离，
并对照你的公式，检验旧的“除以三得千米数”
这一经验规则是否成立。
\end{exercise}

\begin{exercise}[$\star\star\star$]\label{exo:g7:motion-average-speed:12}
一道老题，值得反复玩味：一名司机
以 \qty{30}{km/h} 的\omterm{def:g7:motion-average-speed:formula}{速度}走完行程的前一半\emph{距离}，
并希望全程的\omterm{def:g7:motion-average-speed:average}{平均速度}达到 \qty{60}{km/h}。请用
总路程除以总时间的方法，以 \qty{60}{km} 的行程为例，
说明后一半路程必须在零时间内走完，
即这一愿望不是难以实现，而是根本不可能实现。
\end{exercise}

\section{问题：快递员的星期五}

\begin{problem}[{周末问题——一名骑车快递员，
一个城市的星期五；调度员的图表，
奖金单中的陷阱}]\label{pb:g7:motion-average-speed:1}
调度员将快递员莉娜的星期五记录在
路程–时间图上，并据此结算她的工资。你来核对
这份表单。

\medskip\noindent\textbf{第一部分——上午，从
图像中读取。}
图像显示：一段直线上升，从（$0$ 分钟，\qty{0}{km}）
到（$30$ 分钟，\qty{9}{km}）；水平段持续至 $45$ 分钟；
此后又是一段直线，到（$75$ 分钟，\qty{15}{km}）；水平段持续至 $90$ 分钟。
\begin{enumerate}
  \item 用文字描述上午的经过：各个路段、停留，
        以及停留可能是为了什么。
  \item 第一段路程的\omterm{def:g7:motion-average-speed:formula}{速度}是多少（\omterm{def:g6:measurement-in-science:measuring}{单位}为 \unit{km/h}）？
  \item 第二段骑行路段的\omterm{def:g7:motion-average-speed:formula}{速度}是多少？
  \item 莉娜在整个 $90$ 分钟的上午中，
        包括停留在内的\emph{平均}\omterm{def:g7:motion-average-speed:formula}{速度}是多少？
\end{enumerate}

\medskip\noindent\textbf{第二部分——下午，从
公式中计算。}
\begin{enumerate}[resume]
  \item 下午第一段：行程 \qty{12}{km}，\omterm{def:g7:motion-average-speed:formula}{速度}稳定为
        \qty{24}{km/h}。用了多少分钟？
  \item 第二段：到市区送货，用时 \qty{20}{minutes}，
        \omterm{def:g7:motion-average-speed:formula}{速度}为 \qty{18}{km/h}，行驶了多少千米？
  \item 第三段：骑车返回站点，路程 \qty{10}{km}，
        用时 $25$ 分钟。匀速行驶的\omterm{def:g7:motion-average-speed:formula}{速度}（\omterm{def:g6:measurement-in-science:measuring}{单位}为
        \unit{km/h}）是多少？
  \item 下午总计：将各段的路程和时间分别相加，
        （只算行驶段，不含停歇），
        车轮滚动时保持的\omterm{def:g7:motion-average-speed:average}{平均速度}是多少？
\end{enumerate}

\medskip\noindent\textbf{第三部分——奖金单的陷阱。}
奖金规则为：“整趟行程的\omterm{def:g7:motion-average-speed:average}{平均速度}
高于 \qty{20}{km/h}，即可获得快骑奖金。”
\begin{enumerate}[resume]
  \item 莉娜星期五全天：共行驶 \qty{37}{km}，用时 $3$ \omterm{ex:g2:measuring-time:clock}{小时}
        （包括所有停留）。她能获得奖金吗？
  \item 她的同事马克斯这样声称：“我一半行程以
        \qty{30}{km/h} 行驶，另一半以 \qty{15}{km/h} 行驶——
        平均下来是 $22.5$！”调度员核对后发现：两半
        各为 \qty{12}{km}。请计算马克斯真正的\omterm{def:g7:motion-average-speed:average}{平均速度}，
        并判断他能否获得奖金。
  \item 请向马克斯用时间的语言解释，他的 $22.5$
        错在哪里。
  \item 莉娜为快递员提出一条更公平的奖金规则——
        不因送货停留而受罚的规则。请提出一条
        （图像能告诉你方法），并说明它度量的是
        哪一种\omterm{def:g7:motion-average-speed:formula}{速度}。
\end{enumerate}

\medskip\noindent\textbf{第四部分——调度员的教训。}
\begin{enumerate}[resume]
  \item 星期五回顾：写出调度员的三句话——
        一段直线、一段水平段和一段较陡段在图像上
        各表示什么意思——
        以及关于求\omterm{def:g7:motion-average-speed:average}{平均速度}的
        那一条警告。
\end{enumerate}
\end{problem}
